\documentclass[10pt,a4paper]{amsart}
\usepackage[margin=19mm]{geometry}
\usepackage{amsmath,amssymb,mathtools}
\usepackage[scr=rsfs,cal=euler]{mathalfa}
\usepackage{xfrac}
\usepackage[unicode,hidelinks,bookmarks=false]{hyperref}
\newcommand{\ie}{i.e.\ }
\newcommand{\cf}{cf.\ }
\newcommand{\resp}{respectively\ }
\newcommand{\Subsection}{\subsection}
\providecommand{\nobreakdash}{\nobreak\hbox{-}}
\setlength{\emergencystretch}{2em}
\multlinegap0pt
\usepackage{bm}
\usepackage{bbm}
\usepackage{dsfont}
\usepackage{xspace}
\usepackage{cancel}
\usepackage{mathtools}
\usepackage{amsthm}
\makeatletter
\@ifundefined{thm}{\newtheorem{thm}{Thm}}{}
\@ifundefined{prop}{\newtheorem{prop}[thm]{Proposition}}{}
\makeatother
\title[Roots in the substitution group and in the group of Riordan ]{Roots in the substitution group and in the group of Riordan matrices with ones in the main diagonal}
\author{Jorge Calero-Sanz, Luis Felipe Prieto-Mart\'inez}
\date{Source: arXiv:2509.25972v1 (CC BY 4.0)}
\begin{document}
\maketitle

\noindent\emph{Source and licence.} Roots in the substitution group and in the group of Riordan matrices with ones in the main diagonal. Version arXiv:2509.25972v1 (CC BY 4.0), distributed under the Creative Commons Attribution 4.0 International licence, as recorded in the arXiv licence metadata for this exact version and shown on the abstract page https://arxiv.org/abs/2509.25972v1. Full rights notice: licenses/arxiv-2509.25972-CC-BY-4.0.txt.\par\smallskip

\noindent\textit{Editorial scope.} Counted unit: the Prop whose source label is \texttt{None} in the article (source0.tex lines 647--664, source label \texttt{None}), delivered together with the article's own complete original proof of it (source0.tex lines 666--708, one \texttt{proof} environment of 43 lines). No mathematics is written, repaired or re-derived: statement and proof are the article's own text line for line. Required context retained uncounted: eq.ass (lines 339--341). Every definition, display and label of those lines is retained unchanged. Omitted from this package and not redistributed: 1--338, 342--646, 665--665, 709--1222. Nothing inside a retained excerpt is omitted or altered: every retained body is delivered exactly as the article's lines are written, so no omission receipt of any kind accompanies this package. Citations: the retained excerpts carry no citation command at all, so no bibliography entry of the article is transcribed and no reference is redistributed. Reference adaptation: none was needed and none was made; every reference inside the delivered unit resolves inside it, so no unresolved reference or citation is produced. Printed slips kept literal: no printed word, symbol, hypothesis or number of the counted statement or of its proof was altered. Layout and class: the article's own class is \texttt{elsarticle} with packages including amssymb, babel, graphicx, epstopdf, epsfig, amsfonts, fancyhdr, graphics, amsmath, titlesec, tikz, bm, amsthm, hyperref; the wrapper is the lane's pinned \texttt{amsart} editorial wrapper at 10pt on A4, which restates the theorem-like environments the retained text needs and adds the article's own macro definitions that the retained text needs (none), with extra wrapper packages (bm, bbm, dsfont, xspace, cancel, mathtools). The delivered numbering is the unit's own. Assets: no retained span contains a figure environment, a graphics-inclusion command, a PDF-page inclusion, a table or a TikZ picture; no image file is copied into this package, the package declares no asset dependency and no image file is redistributed with it. Licence: version arXiv:2509.25972v1 of this article is distributed under the Creative Commons Attribution 4.0 International licence, as recorded in the arXiv licence metadata for this exact version and shown on the abstract page https://arxiv.org/abs/2509.25972v1. A full rights notice is delivered at licenses/arxiv-2509.25972-CC-BY-4.0.txt. Characters: every retained excerpt is pure ASCII, so the delivered engine, XeLaTeX with its default Latin Modern text font, reports no missing character.
\begin{equation} \label{eq.ass} \resizebox{.9\textwidth}{!}{$   R(1,\omega) = \left[\begin{array}{lllllllll} 1\\ 0 & 1  \\ 0 &\omega_2 & 1 \\0&\omega_3 & 2\omega_2& 1\\0& \omega_4&2\omega_3+\omega_2^2 & 3\omega_2& 1\\0& \omega_5&2\omega_4+2\omega_2\omega_3 & 3\omega_3+3\omega_2^2& 4\omega_2& 1\\ 0& \omega_6& 2\omega_5+2\omega_2\omega_4+\omega_3^2& 3\omega_4+6\omega_2\omega_3+\omega_2^3& 4\omega_3+6\omega_2^2& 5\omega_2&1 & \\
 0& \omega_7 & 2\omega_6+2\omega_2\omega_5+2\omega_3\omega_4 & 3\omega_5+6\omega_2\omega_4+3\omega_3^2+3\omega_2^2\omega_3 &4\omega_4+12\omega_2\omega_3+4\omega_2^3 &5\omega_3+10\omega_2^2 & 6\omega_2 & 1\\
 \vdots & \vdots&\vdots &\vdots & \vdots& \vdots&  \vdots&\vdots  &  \ddots \end{array}\right] $}\end{equation}

\begin{prop}
    
 Let $g=x+g_2x^2+g_3x^3+\ldots\in\mathcal J(\mathbb Z_2)$, then  $R_{7}(1,g)$ admits a square root $R_{7}(1,\omega)$ in $\mathcal{L}_{7}'(\mathbb Z_2)$  if and only if $g_2=g_3=0$ and we are in one of the following cases: 


\begin{itemize}

\item  $g_4=0$, $g_5=g_7=0$,

\item $g_4=0$, $g_5=g_7=1$,

\item  $g_4=1$, $g_5=g_7=1$.



\end{itemize}

\end{prop}

\begin{proof} 
In view of the matrix in Equation \eqref{eq.ass}, it is clear that the matrices $R_{7}(1,\omega) \in \mathcal L_7'(\mathbb Z_2)$ are of the type
\begin{eqnarray*} 
\underbrace{\left[\begin{array}{l l l l l l l l} 1& & & & & & & \\ 0& 1& & & & & & \\
0&\omega_2 &1 & & & & & \\
0&\omega_3 &0 &1 & & & & \\
0& \omega_4&\omega_2 & \omega_2&1 & & & \\
0&\omega_5 &0 &\omega_2+\omega_3 & 0&1 & & \\
0&\omega_6 & \omega_3&\omega_2+\omega_4 & 0 & \omega_2&1 & \\
0&\omega_7 &0 &\omega_3+\omega_5+\omega_2\omega_3 &0 &\omega_3 & 0&1 \\
 \end{array}\right]}_{R_{7}(1,\omega)} \Rightarrow 
\underbrace{\left[\begin{array}{l l l l l l l l} 
1& & & & & & & \\ 
0& 1& & & & & & \\
0& 0 &1 & & & & & \\
0&0 &0 &1 & & & & \\
0& \omega_{2} + \omega_2 \omega_3& 0 & 0&1 & & & \\
0&\omega_{3} + \omega_2 \omega_3 &0 &0 & 0&1 & & \\
0& \omega_3 \omega_4 + \omega_2 \omega_5 & 0 & \omega_{2} + \omega_2 \omega_3 & 0 & 0&1 & \\
0&\omega_{3} + \omega_2 \omega_{3}&0 & \omega_{3} + \omega_{2}\omega_{3} &0 &0 & 0&1 \\
\end{array}\right]}_{R_{7}(1,\omega)^2}  .\end{eqnarray*} 

\noindent Imposing that the elements in the 1-column of $(R_7(1,\omega))^2$ equals the elements in the 1-column of $R_7(1,g)$, where we have clearly that $g_2 = g_3 = 0$, and we need that:
$$\begin{cases}\omega_2+\omega_2\omega_3=g_4\\ 
\omega_3+\omega_2\omega_3=g_5 \\
\omega_2\omega_5+\omega_3\omega_4=g_6 \\
\omega_3+\omega_2\omega_3=g_7\end{cases} \Longrightarrow \begin{cases}\omega_2+\omega_2\omega_3=g_4\\ 
\omega_3+\omega_2\omega_3=g_5=g_7 \\
\omega_2\omega_5+\omega_3\omega_4=g_6 \end{cases}$$

\noindent Now, let us do a case by case study of this system:
\begin{itemize}
\item $\omega_2=0,\omega_3=0\Rightarrow g_4=g_5=g_6=g_7=0$,

\item $\omega_2=0,\omega_3=1\Rightarrow g_4=0$, $g_5=g_7=1$,

\item $\omega_2=1,\omega_3=0\Rightarrow g_4=1$, $g_5=g_7=0$,

\item $\omega_2=1,\omega_3=1\Rightarrow g_4=g_5=g_7=0$.

\end{itemize}
$ $
\end{proof}

\end{document}
